\documentclass[11pt]{article}
\usepackage[T1]{fontenc}
\usepackage{lmodern,microtype}
\usepackage{amsmath,amssymb}
\usepackage[letterpaper,margin=1in]{geometry}
\usepackage{booktabs,array,enumitem}
\usepackage[colorlinks=true,linkcolor=blue!50!black,urlcolor=blue!50!black]{hyperref}
\newcommand{\GL}{\mathrm{GL}}
\newcommand{\Gb}{\overline G}
\newcommand{\GDb}{\overline{G_D}}
\title{Referee report on \texttt{mjl\_part4\_draft2.tex}\\
\normalsize ``A metaplectic Jacquet--Langlands transfer for division
algebras of prime degree, IV'', Draft~2}
\author{Claude (Anthropic), for David Joyner}
\date{November 7, 2026}
\begin{document}
\maketitle

\section*{0. Scope}
\paragraph{Read in full.} Draft~2 (966 lines). It compiles to 15 pages
without errors.

\paragraph{Checked against the attached sources.} Every statement that
Draft~2 cites from:
\begin{itemize}
\item Part~I, Draft~6;
\item Part~II, Draft~6 (theorem numbers only);
\item Part~III, Draft~6;
\item [FK] (\S\S3, 4, 6, and the Introduction, for $(n,N)=1$);
\item Zou, arXiv v3: (2.2), (3.2), Remark~3.1, Lemma~3.5, and the
definition of $x^*$ on p.~8.
\end{itemize}
For Parts~I--III I relied on my reading for the Draft~1 report and
rechecked only the cited items.

\paragraph{Version caveat.} Draft~2 says it was checked against Part~I,
Draft~5, Part~II, Draft~4, and Part~III, Draft~7. I have Part~III only
in Draft~6. All the Part~III numbers Draft~2 cites agree with Draft~6:
\begin{itemize}
\item Lemma~3.3;
\item Theorems~4.3, 7.6 and~9.1;
\item Corollaries~5.1 and~5.2;
\item Lemma~7.5;
\item Proposition~11.3.
\end{itemize}
If Draft~7 renumbered anything, these must be rechecked.

\paragraph{Numbering used below (compiled).}
\begin{itemize}
\item Proposition~2.1 (diagonal commutator);
\item Theorem~3.1 (FK good elements);
\item Example~3.2 and Remark~3.3;
\item Lemma~4.1 and Proposition~4.2 (kernel, index, modulus);
\item Remark~4.3 and Example~4.4;
\item Proposition~5.1 (Hasse principle) and Theorem~5.2 (global good classes);
\item Proposition~6.1 (lattice criterion);
\item Proposition~7.1 (Zou constant);
\item Proposition~8.1 (conditional comparison);
\item Proposition~9.1.
\end{itemize}

\section*{Summary verdict}
Draft~2 is a large improvement. It deals with every substantive point of
my Draft~1 report:
\begin{itemize}
\item the Flicker--Kazhdan good-element theorem is attributed and labeled
\textbf{[E]};
\item the global good-class criterion is now Theorem~5.2;
\item the two global comparisons are clearly separated;
\item $b_D$ is shown to be independent of the choice of measures;
\item $(n,N)=1$ and Zou are cited;
\item the central characters are fixed in the right order.
\end{itemize}

\textbf{I found no false mathematical statement.} One proof step is
wrong as stated: the general Haar-modulus claim in Proposition~4.2. Its
conclusion is right, but the stated reason is false in general (\S1.3).

The main remaining issues are:
\begin{enumerate}[label=(\arabic*)]
\item Proposition~8.1 is close to a tautology. It should be replaced by a
precise local matching condition at split places, which [FK, \S6]
already supplies (\S1.6).
\item The paragraph on Li's trace formula says something false about real
places and makes an unsupported claim (\S1.7).
\item Drafting history is still in the body (\S3).
\end{enumerate}

\begin{center}\small
\begin{tabular}{@{}p{0.31\linewidth}p{0.17\linewidth}p{0.44\linewidth}@{}}
\toprule
Result & Mark & Comment\\\midrule
Prop.~2.1 & correct & Rechecked; the $c$/$c_v$ issue is resolved\\
\S2.2 (star map) & correct; minor repair & Say why $u(x)=1$\\
Thm.~3.1 \textbf{[E]} & correct & Proof verified; attributed\\
Ex.~3.2, Rem.~3.3 & correct & The regularity remark is valuable\\
Lemma~4.1 & correct & \\
Prop.~4.2 & correct; proof repair & The general ``$[A:I]/|K|$'' rule is false; the formula holds for another reason\\
Prop.~5.1 \textbf{[E]}, Thm.~5.2 & correct & \\
\S\S6.1, 6.3 (example) & correct & I checked $\Delta(t)/\Delta(t^3)$ and the values 9 and 27\\
Prop.~6.1 & correct & \\
Prop.~7.1 & correct; minor repair & ``Provided $x^n$ is regular'' is automatic; justify the central-character hypothesis\\
Prop.~8.1 \textbf{[C]} & correct but nearly empty & Replace (i)--(ii) by the FK~\S6 matching condition\\
\S8.4 (Li) & false statement & $K$ has no real places; ``known'' is uncited\\
Prop.~9.1 & correct & This is Prop.~6.1 again; consider merging\\
\bottomrule
\end{tabular}
\end{center}

\section*{1. Correctness}

\subsection*{1.1 Section 2}
\paragraph{Proposition~2.1.} Correct. It is the computation I verified for
Draft~1. The distinction between $c$ and $c_v$ is now explicit, and the
remark on $(r,n)=1$ is right.

\paragraph{Star map (\S2.2).} Correct, with one justification to add.
\begin{itemize}
\item \emph{The power identity.} The identity
$(\tilde h_1\tilde h_2)^n=\tilde h_1^n\tilde h_2^n[\tilde h_2,\tilde h_1]^{n(n-1)/2}$
for a central commutator, and its vanishing for $n$ odd, are fine.
\item \emph{Why $\zeta^*=1$.} The products $\sigma(\zeta^kI,\zeta I)$ are
powers of $(\zeta,\zeta)_n=(\zeta,-1)_n=1$. I checked this.
\item \emph{Why $u(x)=1$.} ``For odd $n$ the correction $u(x)$ is
identically~$1$'' is true, but Zou's Lemma~3.5 proves only that the
\emph{ratio} $u(x)/u(x_1)u(x_2)$ is~$1$ for odd $n$. The real reason is
on Zou, p.~8: $u(x)=\pm1$ ``as in [FK, \S4]'', and $u(x)$ lies in
$\mu_n$. For odd $n$, $\{\pm1\}\cap\mu_n=\{1\}$. Say this.
\end{itemize}

\subsection*{1.2 Section 3}
\paragraph{Theorem~3.1.} Correct. I checked the proof:
\begin{itemize}
\item the projection formula $(a,N_{E_i/F}b)_F=(a,b)_{E_i}$ for
$a\in F^\times$;
\item the step $a^{r(1+2c_v)}=a\cdot a^{r-1+2rc_v}$.
\end{itemize}

\paragraph{Remark~3.3.} Correct and useful. Two small additions:
\begin{itemize}
\item In the example, $\mathrm{diag}(1,\zeta_3)$ is regular but not good,
since $\zeta_3\notin F^{\times3}$ (because $\zeta_9\notin\mathbb Q_3(\zeta_3)$).
\item The last sentence (field tori preserve regularity) has a one-line
proof that should be included. If $t^n\in F$, then $[F(t):F]$ divides
both $n$ and $r$.
\end{itemize}

\subsection*{1.3 Section 4}
\paragraph{Lemma~4.1.} Correct.

\paragraph{Proposition~4.2, kernel and index.} Correct. I checked:
\begin{itemize}
\item $\ker P_n\cong\mu_n^k/\mu_n$, using $\mu_n(E_i)=\mu_n(F)$;
\item $|M^\times/M^{\times n}|=n^2|n|_M^{-1}$;
\item the quotient computation, using Lemma~4.1;
\item Example~4.4: the values $1,9,9$ and $3,81,27$.
\end{itemize}

\paragraph{Proposition~4.2, the Haar formula: proof repair.} The formula
\[
\int_{A_E}q(a^n)\,da=n^{k-1}|n|_F^{-(r-1)}\int_{P_n(A_E)}q(b)\,db
\]
is correct. The stated reason is not. The proof, and Remark~4.3, assert
that for an open endomorphism $\varphi$ with finite kernel $K$ and image
$I$ of finite index, the density is $[A:I]/|K|$. This fails already for
$A=\mathbb Z$ (counting measure) and $\varphi(a)=na$. The pushforward of
counting measure is counting measure on $n\mathbb Z$, so the density is
$1$, not $n$.

The correct general statement: take a compact open subgroup $U$ with
$U\cap K=1$. Then $\varphi^{-1}(\varphi(U))=UK$, and
\[
\lambda=\frac{\mathrm{vol}(UK)}{\mathrm{vol}(\varphi U)}
=|K|\,[U:\varphi(U)].
\]
For $A_E$, take $U$ a small torsion-free open subgroup of
$\mathcal O_E^\times/\mathcal O_F^\times$. This is a $\mathbb Z_p$-module
of rank $(r-1)[F:\mathbb Q_p]$, so $[U:U^n]=|n|_F^{-(r-1)}$, and
$\lambda=n^{k-1}|n|_F^{-(r-1)}$.

That this equals $[A_E:I]/|K|$ is a coincidence of this group: the
valuation lattice has rank $k-1$ and $|K|=n^{k-1}$. Remark~4.3 should
therefore read: ``the factor is $|\ker P_n|$ times the local index
$[U:U^n]$, not the index $[A_E:P_nA_E]$''.

\subsection*{1.4 Section 5}
\paragraph{Proposition~5.1 and Theorem~5.2.} Correct. They agree with the
corollary I proposed, and the proof is complete. Two remarks:
\begin{itemize}
\item The $D$-side case $v\notin S$ goes through Part~I, Proposition~2.6
(the isomorphism $\GDb(K_v)\cong\Gb_v$). Cite it by number, not as ``the
exact Part~I cover isomorphism''.
\item Remark~5.3 reproduces Part~I's $c(\gamma)$ exactly (Part~I,
\S6.1). Its last sentence is an important point, made correctly.
\end{itemize}

\subsection*{1.5 Sections 6 and 7}
\paragraph{\S6.1.} The formula for $L_D$ agrees with Part~III,
Definition~4.1. The formula for $\Delta$ is correct: the eigenvalues of
$\mathrm{Ad}(t)$ are $\lambda_i/\lambda_j$. The statement about the
identity value is correct.

\paragraph{\S6.3, the explicit example.} Correct. I checked:
\begin{itemize}
\item $\Delta(t)/\Delta(t^3)=|\lambda_1\lambda_2|/|\lambda_1^2+\lambda_1\lambda_2+\lambda_2^2|$;
\item $\lambda_1^2+\lambda_1\lambda_2+\lambda_2^2=\mathrm{Trd}^2-\mathrm{Nrd}$;
\item the limit $9$ at scalars, giving the value $27$.
\end{itemize}
The paragraph on twisted finite quotients corrects my own Draft~1
suggestion, which was sloppy: $\tilde H/\tilde ZU$ loses the genuine
central character.

\paragraph{Proposition~6.1.} Correct, with the redundancy removed.

\paragraph{Proposition~7.1.} Correct; I checked
$1/(d_F(n)|n|_F^{r/2})=b_D/n$. Two repairs:
\begin{itemize}
\item \emph{A redundant proviso.} ``Provided $x^n$ is regular'' is
automatic for a field centralizer, by Remark~3.3. State it as a
consequence.
\item \emph{The central-character hypothesis.} Zou, Remark~3.1, says Zou
deliberately removes the central-character restriction. So the
hypothesis $\omega_\pi|_{\mu_n}=1$ is a genuine restriction, not a
convention. Note what happens without it:
\begin{itemize}
\item if $\omega_\pi|_{\mu_n}\ne1$, the root sum
$\sum_\zeta\omega_\pi(\zeta)$ vanishes, so (3.2) forces
$\theta_{\tilde\pi}(x^*)=0$ at elliptic $x$;
\item for a lift pair, the compatibility
$\omega_\pi(a)=\omega_{\tilde\pi}(a^*)$, together with $\zeta^*=1$,
gives $\omega_\pi|_{\mu_n}=1$ automatically.
\end{itemize}
Saying this would make the proposition self-contained.
\end{itemize}

\subsection*{1.6 Section 8: replace Proposition~8.1}
Proposition~8.1 is correct but carries little content. Hypothesis~(ii)
is essentially the conclusion. More importantly, it leaves the
split-place question vague (``after any prescribed local root
summations''). The question has a definite answer.

\paragraph{Why a root sum cannot be the local identity at split places.}
The $D^\times$-side sum runs over rational classes $[s]$. The cover side
runs over $[s^n]$, in bijection with them by Theorem~5.2. At $v\notin S$,
the local roots $\zeta s$, with $\zeta\in\ker P_n$ (Proposition~4.2), are
not rational. So a local identity of the root-sum form
\[
J(s^n,\tilde f_v)=\lambda_v\sum_\zeta J(\zeta s,f_v)
\]
cannot be used term by term. What is needed is
\[
J(s^n,\tilde f_v)=\lambda_v\,J(s,f_v)\quad\text{for all regular }s.
\]
This forces $J(\cdot,f_v)$ to be constant on cosets of $\ker P_n$.

\paragraph{FK already impose exactly this.} [FK, \S6] state that they
``deal only with functions $f$ with the property
$\omega(z)F(x,f)=F(x',f)$ for any $x,x'$ with $zx^*=x'^*$''.
(I read this in the attached scan.) Proposition~8.1(i)--(ii) should be
replaced by:
\begin{enumerate}[label=(\alph*)]
\item at each $v\notin S$, the ordinary test function satisfies the FK~\S6
invariance;
\item a pointwise local identity holds with an explicit constant
$\lambda_v$;
\item the product of all the $\lambda_v$, the values $b_{D,v}$ at $v\in S$,
and the archimedean factors equals~$1$.
\end{enumerate}

\paragraph{A second constraint, from Remark~3.3.} Take $s$ regular with
$s^n$ singular (e.g.\ $s=\mathrm{diag}(1,\zeta)$). The map $s\mapsto s^n$
is a local homeomorphism of the torus. So $J(s^n,\tilde f_v)$, pulled
back to a neighborhood of $s$, must be locally constant there, because
$J(\cdot,f_v)$ is locally constant on the regular set. But near a
central $s^n$, $J(\cdot,\tilde f_v)$ has a non-constant germ expansion.
So the existence of the matching $f_v$ imposes vanishing conditions on
germs of $\tilde f_v$. Alternatively it forces a support condition, as in
Zou's restriction to $x$ with $x^n$ regular.

This should be written into the ledger. It is a precise, checkable
obligation, and it shows that matching at split places is not merely a
pullback.

\paragraph{The rest of \S8.} \S8.1--\S8.2 are correct:
\begin{itemize}
\item the product formula $\prod_v\Delta_v(\gamma)=1$;
\item the local identity at $v\in S$ (Part~III, Theorem~9.1);
\item the formal-degree conversion $d(\mathrm{St})\,dg_{D,v}$.
\end{itemize}
The citation ``Part~II, Theorem~6.2'' is right: \texttt{thm:comparison2}
is 6.2, and Theorem~C$'$ is 6.5. Delete the phrase ``in the numbering
cited by the referee''.

\subsection*{1.7 \S8.4: Li's trace formula}
``Under the trace Paley--Wiener hypothesis at real archimedean places
\ldots\ This hypothesis is known for the cited metaplectic covers of
$\GL$'' needs correcting:
\begin{itemize}
\item $K$ is totally imaginary, so there are no real places.
\item Li's hypothesis [Li4] concerns the invariant Paley--Wiener theorem at
the archimedean places.
\item At complex places the covers split (Part~I), so the hypothesis
reduces to the linear case. That case is a known theorem, but it needs a
citation (e.g.\ Clozel--Delorme), which the draft does not give.
\item ``Known for the cited metaplectic covers'' is unsupported as it
stands.
\end{itemize}
Part~II, \S7(iii), has the right wording.

\subsection*{1.8 Section 9}
\paragraph{\S9.1.} Correct.
\begin{itemize}
\item Note that $\theta(a)=\tilde\theta(a^*)$ is automatically trivial on
$K^\times$, since $a^*=\iota(a)^n=\iota(a^n)$ and $\tilde\theta$ is
trivial on $\iota(K^\times)$. Worth one sentence.
\item Proposition~9.1 restates Proposition~6.1. Its \textbf{[C]} label
is misleading: the implication itself is proved. Merge it into
Proposition~6.1 as a remark.
\end{itemize}

\section*{2. Consistency and citations}
\begin{enumerate}[label=(\arabic*)]
\item \emph{Labels.} \textbf{[R]} is gone. The meanings of
\textbf{[O]} and \textbf{[C]} still differ from Part~I's (``proof
outlined'', ``conjectural''). This needs a series-wide decision.
\item \emph{Titles.} The title now follows Parts~I--II. But the
bibliography entry [JoynerIII] still gives Part~III's old ``Toward
\ldots'' title. The author line (Joyner, ChatGPT, Claude) differs from
Parts~I--II. ``D.~Joyner with collaborators'' in the bibliography should
give the actual authors.
\item \emph{Citation style.} Only DKV, Kottwitz, Li and Zou are cited with
\verb|\cite|. FK, KP and Parts~I--III are typed in by hand. [Takeda] is
in the bibliography but never cited.
\item \emph{Unverified.} Not checked by me:
\begin{itemize}
\item Takeda's page range: Draft~2 gives 415--449, Draft~1 gave none;
I could not check it;
\item KP (Publ.\ Math.\ IH\'ES~59, 35--142), which agrees with the
standard citation as far as I know;
\item the DKV page range;
\item Kottwitz and Li, which I did not have.
\end{itemize}
\item \emph{Checked.} The FK and Zou items listed in \S0; Part~I's
Proposition~2.3(b), Proposition~3.1, Proposition~3.5(a), \S6.1 and
$c(\gamma)$; Part~II's Definition~2.1 and theorem numbers.
\end{enumerate}

\section*{3. Readability}
\begin{enumerate}[label=(R\arabic*)]
\item \emph{Drafting history is still in the body.} Move these to the log:
\begin{itemize}
\item the paragraph ``Scope of direct cross-checking'' (\S1), which
lists drafts and mentions ``Claude's attached referee report'';
\item in \S2.1, ``this is exactly the cocycle used in Part~IV, Draft~1;
its letter $c$ meant \ldots'';
\item in \S8, ``in the numbering cited by the referee''.
\end{itemize}
The acknowledgments may credit the referee report, as Parts~I--II do.
``Suggested repairs to the chronology of the series'' is unclear and
should be reworded or cut.
\item \emph{Disclaimers.} These are reduced, but they still recur:
\begin{itemize}
\item the end of the abstract;
\item the status convention (``No representation-theoretic correspondence
follows \ldots'');
\item Remark~7.2;
\item \S8 (``not a formal consequence'');
\item the last paragraph of \S9;
\item the acknowledgments.
\end{itemize}
The status table suffices.
\item \emph{Orienting paragraphs.} These are good throughout.
\item \emph{Glossary.} It is now short and specific to the paper. That is
acceptable.
\item \emph{Abstract.} It is accurate. ``Shows why regularity of a root must
be distinguished from regularity of its $n$th power'' refers to
Remark~3.3, not to Proposition~4.2. Reword.
\end{enumerate}

\section*{4. Recommendation}
Minor revision for \S\S2--7, after these changes:
\begin{itemize}
\item repair the proof of Proposition~4.2 and Remark~4.3;
\item justify $u(x)=1$;
\item make the hypothesis of Proposition~7.1 explicit.
\end{itemize}

For \S8, the most valuable next step is to replace Proposition~8.1 by
the FK~\S6 matching condition at split places, together with the germ
constraint of \S1.6. This would turn \S8 into a precise statement of the
remaining local theorem. Better still, the theorem could be proved:
\begin{itemize}
\item for $v\nmid n$, via FK's unit-element and matching results under
$(n,N)=1$;
\item at $v\in S$, using Part~III, Theorem~9.1.
\end{itemize}

The explicit wild test of \S6.3 remains the cleanest way to settle the
first case of the integrality question. It can be done as a computation
in a twisted group algebra in GAP or Sage.

\end{document}
